\par\textbf{Procedimiento completo para comparar tu solución.} La estimación de base de nube usa diferencia de temperaturas, para la que grados Celsius y kelvin tienen el mismo incremento. A 1 km, se aplica por separado el gradiente de la parcela seca y el ambiental.

\begin{enumerate}

\item Aplicar la correlación docente suministrada.

\[\fitmath{\mathrm{LCL}=125\cdot \left(23-11\right)\approx 1500\,\mathrm{m}}\]

\item Temperatura de la parcela a 1 km.

\[\fitmath{\mathrm{Tparcel}=23-9.8\approx 13.2\,\mathrm{degC}}\]

\item Temperatura ambiental a la misma altura.

\[\fitmath{\mathrm{Tenv}=23-6.5\approx 16.5\,\mathrm{degC}}\]

\item Comparar parcela con ambiente; un valor negativo indica parcela más fría en este modelo.

\[\fitmath{\mathrm{difference}=23-9.8-\left(23-6.5\right)\approx -3.3\,\mathrm{K}}\]

\end{enumerate}
